\subsection*{Problema 17}

\textbf{Enunciado:} Calcular la integral:
\begin{align*}
\int \dfrac{dx}{5 - 4\cos x}
\end{align*}

\textbf{Solución:} 

Para resolver esta integral trigonométrica, 
utilizamos la sustitución del ángulo medio, 
también conocida como sustitución de Weierstrass. 
Definimos la variable auxiliar:
\begin{align*}
z = \tan\left(\dfrac{x}{2}\right)
\end{align*}

A partir de esta sustitución, 
las identidades trigonométricas 
fundamentales se expresan como:
\begin{align*}
\cos x &= \dfrac{1 - z^2}{1 + z^2}, \quad
dx = \dfrac{2\,dz}{1 + z^2}
\end{align*}

Sustituimos estas expresiones 
dentro de la integral original:
\begin{align*}
\int \dfrac{\dfrac{2\,dz}{1 + z^2}}
{5 - 4\left(\dfrac{1 - z^2}{1 + z^2}\right)}
\end{align*}

Simplificamos el denominador 
multiplicando ambos términos 
por $(1 + z^2)$:
\begin{align*}
&= \int \dfrac{2\,dz}
{5(1 + z^2) - 4(1 - z^2)} \\
&= \int \dfrac{2\,dz}
{5 + 5z^2 - 4 + 4z^2}
\end{align*}

Agrupamos los términos semejantes 
en el denominador:
\begin{align*}
&= \int \dfrac{2\,dz}{1 + 9z^2}
\end{align*}

Factorizamos el coeficiente $9$ 
del denominador para llevar la integral 
a la forma estándar del arco tangente:
\begin{align*}
&= \int \dfrac{2\,dz}
{9\left(\dfrac{1}{9} + z^2\right)} \\
&= \dfrac{2}{9} \int \dfrac{dz}
{\left(\dfrac{1}{3}\right)^2 + z^2}
\end{align*}

Aplicamos la fórmula de integración estándar:
\begin{align*}
\int \dfrac{du}{a^2 + u^2} 
= \dfrac{1}{a} \arctan\left(\dfrac{u}{a}\right) + C
\end{align*}

En nuestro caso, $a = \dfrac{1}{3}$ y $u = z$:
\begin{align*}
&= \dfrac{2}{9} \left( \dfrac{1}{\dfrac{1}{3}} 
\arctan\left(\dfrac{z}{\dfrac{1}{3}}\right) \right) + C \\
&= \dfrac{2}{9} \left( 3 \arctan(3z) \right) + C \\
&= \dfrac{2}{3} \arctan(3z) + C
\end{align*}

Finalmente, regresamos a la variable 
original $x$ sustituyendo $z = \tan\left(\dfrac{x}{2}\right)$:
\begin{align*}
\int \dfrac{dx}{5 - 4\cos x} 
= \dfrac{2}{3} \arctan\left(3 \tan\left(\dfrac{x}{2}\right)\right) + C
\end{align*}

\begin{figure}[H]
	\centering
	\includegraphics[width=\columnwidth]{semana_04/graficas/SEM04_P17_grafica.png}
	\caption{Representación gráfica de $f(x)$ y su antiderivada $F(x)$ evaluada para la constante de integración $C = 0$ (Problema 17).}
	\label{fig:problema17}
\end{figure}
